\clearpage
\fancyhead[L]{\small CONNAISSANCE DES TEMS.}
\hypertarget{native-559-lower}{}\label{native-559-lower}
\begin{center}\small Göttingische gelehrte Anzeigen. Stück 8. Seite 873..876. 3. Juni 1813.\\[3pt]\rule{.72\linewidth}{.3pt}\end{center}

\emph{Connaissance des tems ou des mouvemens célestes, à l'usage des astronomes et des navigateurs, pour l'an} 1813, publiée par le bureau des longitudes. Juillet 1811. Paris. 228 Seiten im Oktavformat.

\emph{Connaissance des tems etc.} pour l'an 1814. Avril 1812. 240 Seiten.

Die Tabelle zur geografischen Lage der wichtigsten Orte der Erde ist erneut erheblich umfangreicher geworden, besonders durch die Verwendung der Bestimmungen, die auf den Reisen von \textsc{Dentrecasteaux}, \textsc{Humboldt} und \textsc{Krusenstern} vorgenommen wurden. So enthält der Jahrgang für 1814 über 1600 Einträge und folglich über 100 mehr als der Jahrgang für 1812. Die \emph{Zusätze} zu diesen beiden Jahrgängen sind diesmal nicht zahlreich. Der Jahrgang 1813 enthält nur zwei kurze Aufsätze von \textsc{Laplace}. Im ersten erweitert er die Theorie der sogenannten Methode der kleinsten Quadrate in einem wichtigen Punkt. Mit einer kunstvollen Analyse zeigt er, dass bei der von dieser Methode vorgeschriebenen Kombination der Grundgleichungen der zu befürchtende mittlere Fehler in den Ergebnissen minimal wird, sobald die Anzahl der Grundgleichungen sehr groß ist, wie auch immer das Wahrscheinlichkeitsgesetz der Fehler lautet. Der Verfasser führt diese Untersuchung in seinem kürzlich erschienenen größeren Werk über die Wahrscheinlichkeitsrechnung ausführlicher aus. Der andere Zusatz betrifft die Mondgleichung mit langer Periode. Auf ihre Notwendigkeit hat man aus der Abweichung zwischen den für 1692, 1756, 1779 und 1801 ermittelten Epochen der mittleren Mondlänge \SourceLeaf{560}geschlossen, ohne sie bisher in der Theorie ganz zufriedenstellend nachweisen zu können. In den Mondtafeln von \textsc{Bürg} hatte man sie proportional zum Produkt des Sinus der Größe aus der doppelten Länge des Mondknotens plus der Länge des Mondperigäums minus der dreifachen Länge des Sonnenperigäums angesetzt und ihren Koeffizienten mit $= -14''$ angenommen. \textsc{Laplace} gibt nun bekannt, dass er gewichtigere Gründe dafür habe (über die er sich aber nicht näher äußert), sie proportional zum Kosinus der Summe aus der doppelten Länge des Mondknotens und der Länge des Mondperigäums anzusetzen, und bestimmt ihren Koeffizienten mit $= -13''96$. (\textsc{Burckhardt} hat in seinen neuen Mondtafeln, wie wir bereits bei deren Besprechung bemerkten. dieselbe Form zugrunde gelegt und den Koeffizienten mit $= -12''5$ angenommen.) Wir sehen der ausführlicheren Darstellung der theoretischen Untersuchungen des Verfassers zu dieser heiklen Frage mit Spannung entgegen.

Die Zusätze zum Jahrgang 1814 sind folgende: 1) Über den Ursprung der Kometen, von \textsc{Lagrange}. Der große Geometer (den der Tod inzwischen, wenn auch in hohem Alter, so doch immer noch zu früh den Wissenschaften entrissen hat) dehnt hier die bekannte Hypothese von \textsc{Olbers} über den Ursprung der neuen Planeten auf die Kometen aus. Mit der Eleganz, die man von ihm stets gewohnt war, entwickelt er die Bedingungen, unter denen ein von einem Planeten auf einer Kreisbahn fortgeschleudertes Bruchstück eine parabolische Bahn beschreibt. Wie man sich leicht überzeugt (\textsc{Lagrange} hat es allerdings nicht ausdrücklich gesagt), sind die Grenzen der relativen Geschwindigkeit, mit der der entstehende Komet fortgestoßen wird, $\BareRoot{2}-1$ und $\BareRoot{2}+1$, wobei die Geschwindigkeit des Planeten als Einheit gilt. Bemerkenswert ist aber der hier von \textsc{Lagrange} aufgestellte (ebenfalls leicht zu beweisende) Satz, dass die Geschwindigkeit $\BareRoot{3}$ die Grenze zwischen rechtläufigen und rückläufigen Kometen bildet, sodass ein (in Bezug auf die Bahn des Planeten) rechtläufiger Komet eine kleinere, ein rückläufiger eine größere Geschwindigkeit erfordert. Man stelle sich einen Planeten in einer hundertmal so großen Entfernung wie unsere Erde vor, der durch die Wirkung eines in seinem Inneren plötzlich, aus welcher Ursache auch immer, frei werdenden elastischen Fluids in mehrere Stücke zersprengt würde. Die Explosion müsste nur so stark sein, dass sie eine zwölf- oder fünfzehnmal so große Geschwindigkeit wie die einer 24-pfündigen Kanonenkugel erteilen könnte, um elliptische oder parabolische Kometen in allen möglichen Dimensionen und in \emph{allen möglichen Richtungen hervorzubringen}. Dieser letzten Behauptung müssen wir jedoch widersprechen. Auf diese Weise könnten nur solche Kometen entstehen, deren Knotenlinie auf der früheren Planetenbahn fast mit der Apsidenlinie zusammenfällt, wenn sie in Sonnennähe bis in den Bereich der Erdbahn herabkommen sollen. Unter sämtlichen bisher berechneten Kometenbahnen ist jedoch nur ein kleiner Teil von der Art, dass man die Möglichkeit einer solchen Entstehung zugestehen könnte. Praktisch arbeitende Astronomen werden überhaupt schwerlich geneigt sein, in Kometen, deren gesamtes Aussehen auf eine völlig andere physische Beschaffenheit hindeutet, Bruchstücke fester Planetenkörper zu erkennen. 2) Einige Notizen aus \textsc{Krusenstern}s Reise um die Welt. 3) Eine sehr praktisch angeordnete Tabelle zur Umrechnung der Sternzeit in mittlere Sonnenzeit, von \textsc{Burckhardt}. 4) Auszug aus den meteorologischen Beobachtungen in Paris im Jahr 1809. Am Ende beider Jahrgänge steht wie gewöhnlich das Verzeichnis der Mitglieder des Längenbüros.
